% !TEX program = xelatex
% =============================================================================
% 中文考研复试简历 LaTeX 模板
% Author: Kody
% Repository: https://github.com/kodyyu1126/Chinese-resume-template-postgraduate
% Compiler: XeLaTeX
% =============================================================================
\documentclass[11pt,a4paper,fontset=none]{ctexart}

% ── 字体设置 ──────────────────────────────────────────────
% 若需在 Overleaf 使用自定义字体，可将已授权的宋体文件放入 fonts/ 文件夹。
% macOS 本地优先使用 Songti SC，英文字母统一使用 Times New Roman。
\IfFontExistsTF{Times New Roman}{\setmainfont{Times New Roman}}{}

\newif\ifuseuploadedcjkfonts
\IfFileExists{fonts/STSong.ttf}{
    \IfFileExists{fonts/STHeiti.ttf}{
        \IfFileExists{fonts/STKaiti.ttf}{\useuploadedcjkfontstrue}{}
    }{}
}{}

\newif\ifusesystemcjkfonts
\IfFontExistsTF{Songti SC}{\usesystemcjkfontstrue}{}

\ifuseuploadedcjkfonts
    \setCJKmainfont[
        Path = fonts/,
        BoldFont = STSong.ttf,
        ItalicFont = STSong.ttf
    ]{STSong.ttf}
    \setCJKsansfont[Path = fonts/]{STHeiti.ttf}
    \setCJKmonofont[Path = fonts/]{STHeiti.ttf}
\else\ifusesystemcjkfonts
    \setCJKmainfont[
        BoldFont = {Songti SC Bold},
        ItalicFont = {Songti SC}
    ]{Songti SC}
    \setCJKsansfont{Songti SC}
    \setCJKmonofont{Songti SC}
\else
    \setCJKmainfont[
        BoldFont = Noto Sans CJK SC,
        ItalicFont = Noto Serif CJK SC
    ]{Noto Serif CJK SC}
    \setCJKsansfont{Noto Sans CJK SC}
    \setCJKmonofont{Noto Sans CJK SC}
\fi\fi

% ── 页面布局 ──────────────────────────────────────────────
\usepackage{geometry}
\geometry{a4paper, left=1.5cm, right=1.5cm, top=1.5cm, bottom=1.5cm}

% ── 颜色定义 ──────────────────────────────────────────────
\usepackage{xcolor}
\definecolor{sectioncolor}{RGB}{46,116,181}
\definecolor{linkcolor}{RGB}{0,112,192}
\definecolor{textcolor}{RGB}{0,0,0}

% ── 标题格式 ──────────────────────────────────────────────
\usepackage{titlesec}
\titleformat{\section}
    {\Large\bfseries\color{sectioncolor}}
    {}
    {0em}
    {\uppercase}
    [\vspace{-0.8em}\rule{\textwidth}{1pt}]
\titlespacing{\section}{0pt}{0.25em}{0.2em}

% ── 列表格式 ──────────────────────────────────────────────
\usepackage{enumitem}
\setlist[itemize]{leftmargin=*, nosep, topsep=0.2em, itemsep=0.1em, parsep=0em}

% ── 图片与排版 ────────────────────────────────────────────
\usepackage{graphicx}
\setkeys{Gin}{draft=false}
\usepackage[export]{adjustbox}
\usepackage{calc}
\usepackage{tikz}

% ── 表格 ─────────────────────────────────────────────────
\usepackage{tabularx}
\usepackage{multirow}

% ── 图标与超链接 ──────────────────────────────────────────
\usepackage{fontawesome5}
\usepackage{hyperref}
\hypersetup{colorlinks=true, linkcolor=linkcolor, urlcolor=linkcolor}

% ── 全局样式 ──────────────────────────────────────────────
\color{textcolor}
\linespread{1.5}
\pagestyle{empty}

% ── 自定义命令 ────────────────────────────────────────────
\newcommand{\daterange}[1]{{\textbf{#1}}}
\newcommand{\entrytitle}[1]{{\textbf{#1}}}
% 经历归属标题使用可嵌入宋体粗体并略微放大，保证跨设备显示与层级区分。
\newcommand{\orgtitle}[1]{{\bfseries\fontsize{12pt}{14pt}\selectfont#1}}
% 项目小标题使用轻微倾斜的宋体，和经历归属标题形成层级。
\newcommand{\projecttitle}[1]{{\IfFontExistsTF{Songti SC}{\CJKfontspec[FakeSlant=0.15]{Songti SC}\itshape}{\itshape}#1}}

% ============================================================
\begin{document}
% ============================================================

% ── 个人信息头部 ──────────────────────────────────────────
\noindent
\begin{minipage}[t]{\textwidth}
    \vspace*{-2.5em}
    \begin{center}
        {\Huge\bfseries\color{sectioncolor} 葛飞翔} \\[0.9em]
        \textbf{电话：}19857950790\hspace{1.5em}
        \textbf{邮箱：}\href{mailto:marshall_gefxiang@163.com}{marshall\_gefxiang@163.com} \\[0.45em]
        \textbf{研究兴趣：} 大模型；异构软硬件推理系统
    \end{center}
\end{minipage}

\vspace{0.3em}

% ── 教育经历 ──────────────────────────────────────────────
\section*{教育经历}
\vspace{-0.3em}

\noindent
\entrytitle{中国科学院大学｜推免硕士研究生 | 人工智能}
\hfill\daterange{2025.09 \textasciitilde{} 至今}

\par\noindent
$\bullet$ \textbf{荣誉经历：} 国科大杭州高等研究院新生奖学金

\noindent
\entrytitle{上海对外经贸大学｜本科 | 数据科学与大数据技术}
\hfill\daterange{2021.09 \textasciitilde{} 2025.06}

\par\noindent
$\bullet$ \textbf{荣誉经历：} 上海市奖学金、爱建特种基金会奖学金、优秀毕业生、优秀共青团员

% ── 业界工作 ──────────────────────────────────────────────
\section*{业界工作}
\vspace{-0.1em}

\noindent
\orgtitle{百度｜模型适配与推理性能优化 ｜ vLLM、SGLang、昆仑芯}
\hfill\daterange{2026.02 \textasciitilde{} 2026.08}

\par\noindent
$\bullet$ \projecttitle{\textbf{ LLM 推理优化：}} 1) 对于Dsv3.2模型，面向昆仑芯 P800，基于 SGLang 对 NSA Indexer、MTP/投机解码、上下文并行开展 profiling、框架对齐与关键函数优化，使单机性能较原准出提升 20\%、PD 分离提升 50\%，并定位解决投机解码乱码等一系列问题；2) 基于上游vLLM Workspace Manager机制 优化vLLM-Kunlun显存管理，使 Qwen3-97B 单机最大上下文翻倍，改动合入 \href{https://github.com/baidu/vLLM-Kunlun/pull/283/changes}{PR 283}；3) 分析算子计算强度，发现 speculative attention 的 FP16 实现较 BF16 快 1.8 倍，替换为更适配硬件的算子实现。

\par\noindent
$\bullet$ \projecttitle{\textbf{框架迁移、开源共建与 AI 提效：}} 1）从零迁移vLLM-Omni到昆仑芯，主导多模态模型在昆仑芯的适配，主流模型的性能达到H20的80\%；2）参与 vLLM-Kunlun 开源共建，贡献者排名第 6/39、合入 9 个 PR，覆盖代码重构、算子优化与显存管理；3）迭代 LLM 部署 SOP Skill 和准出 Agent，自动化模型部署、性能/精度测试及结构化报告。


\par\noindent
$\bullet$ \projecttitle{\textbf{多模型适配与交付：}}完成 Qwen3/Qwen3.5、Wan、Qwen-Image、GLM-Image、Qwen-Image-Edit 的推理适配、性能精度验证与交付，对于小模型诸如 Grounding SAM、OpenFold/SAM3、BEVFormer 等完成训练或推理场景适配，并解决过程中暴露的一系列问题。


\par\noindent
$\bullet$ \projecttitle{\textbf{Dynamo 与 HiCache 特性适配：}}完成 SGLang + HiCache + Dynamo 整体框架调通适配及 GLM 系列验证；修复 KV-aware 模式 DP rank 不一致 Hang 和前缀复用调度失效问题，并在 Dynamo/Router 注入 Prometheus 指标拆解链路。经分析和优化，Dynamo 可稳定提升 TPOT 10\%+，HiCache 在特定数据集使前缀命中率翻倍并显著降低 TTFT。

\noindent
\orgtitle{eBay 软件工程有限公司｜AI Platform工程师 | Ray、Kubernetes、集群调度}
\hfill\daterange{2024.10 \textasciitilde{} 2025.08}

\par\noindent
$\bullet$ \projecttitle{\textbf{平台工程：}}深入实践 Ray Data、Ray Serve 等组件，解决多类型资源标签调度和 PyArrow 内存泄漏问题；相关修复已合入公司内部 Ray 分支（\href{https://github.com/ray-project/ray/pull/49124}{PR 49124}、\href{https://github.com/ray-project/ray/pull/49756}{PR 49756}）。参与内部AI工作流搭建工具开发，主要负责可观测性板块和部分核心功能开发。
\par\noindent
$\bullet$ \projecttitle{\textbf{集群优化与技术跟进：}}负责测试业务部门模型的GPU集群利用情况，尽可能优化资源调度、内存管理、分布式策略等来提升GPU集群的利用率。曾修复 KubeRay 资源调度缺陷使 GPU 利用率翻倍，并探索 LLM 强化学习训练的资源需求、调研框架。




\newpage
% ── 项目经历 ──────────────────────────────────────────────
\section*{项目经历}
\vspace{-0.1em}

\noindent
\orgtitle{IEEE AICAS 2026 - 面向AI芯片的VLM高效推理与优化赛道(rank 6 / 514)}

\par\noindent
$\bullet$ \projecttitle{\textbf{端到端推理加速：}}面向 PPU-ZW810E 优化 Qwen3-VL-2B-Instruct，在 VQA 单 batch、1K 输入 / 128 token 输出评测下，将 TTFT 从 vLLM 的约 77\,ms 降至约 17\,ms，吞吐从约 220 提升至约 340 tok/s；综合精度、TTFT 与吞吐指标取得赛道第 6 / 514 名。

\par\noindent
$\bullet$ \projecttitle{\textbf{全栈优化设计：}}构建图像、文本Radix Tree 等多层级KV Cache，结合 \texttt{torch.compile}、Prefill/Decode CUDA Graph 多策略录图进行推理加速；以 Triton、CUDA 和 TileLang 实现 QK RMSNorm-RoPE、FlashAttention/FlashDecode、FFN、KV 写入等关键融合算子，并通过 trace 驱动的代码优化减少 D2H 同步与小算子 launch。

\noindent
\orgtitle{基于 RISC-V AI CPU 的推理优化}

\par\noindent
$\bullet$ \projecttitle{\textbf{硬件感知 GEMM 优化：}}针对 RISC-V AI CPU 上 Qwen3.5-2B 的 GEMM，结合 Roofline 与 PMU 定位 VPU 发射吞吐和非矩阵指令瓶颈；重排为 N-major 循环、将 Q4\_1 切换至 Q4\_0、把量化 scale 分组由 32 扩至 64，精简汇编指令。在 128-token 场景下，Prefill 相比基线提升 289\%（3.89$\times$），Decode 提升 57\%（1.57$\times$）。

\par\noindent
$\bullet$ \projecttitle{\textbf{模型压缩与数据布局优化：}}将 Vision Encoder 从 24 个 ViT block 精简至 7 个，在精度损失可控前提下使 视觉 prefill 阶段吞吐翻倍； 针对 SSM \texttt{build\_conv\_state} 的非连续访存，设计 64$\times$64 分块转置与 L1 Cache 复用，实现单测 7--10$\times$、端到端约 10\% 加速。


% ── 研究工作 ──────────────────────────────────────────────
\section*{研究工作}
\vspace{-0.1em}

\par\noindent
$\bullet$ \href{https://openreview.net/pdf?id=cs7cFJRYLj}{\entrytitle{Dynamical Adapter Fusion: Constructing A Global Adapter for Pre-Trained Model-based Class-Incremental Learning}}，\entrytitle{RuiQi Liu
, Boyu Diao, Zijia An, Runjie Shao, Feixiang Ge, et al. NeurIPS 2026, CCF-A。}

\noindent
$\bullet$ \href{https://doi.org/10.1109/IJCNN64981.2025.11227768}{\entrytitle{Using Domain Adapter and Prompt Learner for Few-shot Teaching Action Recognition in videos}}，\entrytitle{Feixiang Ge, Xinlei Li, Yan Rong, et al.，IJCNN 2025；第一作者，CCF-C。}

\par\noindent
$\bullet$ \href{https://doi.org/10.1007/978-981-96-9891-2_39}{\entrytitle{A Hybrid CNN-Transformer Network for Fine-Grained Tongue Multi-Attribute Classification with Tongue Group Attribute Mask Training in Traditional Chinese Medicine Diagnosis}}，\entrytitle{Liu J, Ge F, Li F, et al.，International Conference on Intelligent Computing，Springer Nature Singapore，2025：467--479；共同一作，CCF-C。}


\par\noindent
$\bullet$  \entrytitle{《通用深度学习模型搭建与开发辅助工具》, 登记号：2026SR0303020；第一著作权人，软件著作权。}


\end{document}
